% ============================================================
%  lang/en.tex — every user-visible string in the English edition
%
%  Nothing in programs/en/ may hard-code a word that a macro also typesets.
%  If a label is needed and is not here, add it here AND in lang/pl.tex; CI
%  fails the build when the two files do not define the same macro set.
% ============================================================

% ---------- Structure ----------
\newcommand{\lblProgram}{Program}
\newcommand{\lblPrograms}{Programs}
\newcommand{\lblFrame}{Frame}
\newcommand{\lblFrames}{Frames}
\newcommand{\lblOutcomes}{What will I learn from this program?}
\newcommand{\lblOutcomesLead}{When you have finished this program you will be able to:}
\newcommand{\lblQuiz}{Quiz}
\newcommand{\lblSummary}{Summary}
\newcommand{\lblCanYou}{Can you?}
\newcommand{\lblOnEntryYouCould}{On entry this program promised you would be able to}
\newcommand{\lblTestExercises}{Test exercises}
\newcommand{\lblFurther}{Further problems}
\newcommand{\lblAnswers}{Answers}
% The opener's frame-range box. Two arguments rather than a bare word, so
% the translator owns the word order and not just the preposition -- but the
% word itself comes from \lblRouteTo, which the Quiz's route boxes also use,
% so the book says to/do from one place and the two cannot drift.
\newcommand{\lblFrameRange}[2]{\lblFrames~#1 \lblRouteTo\ #2}
\newcommand{\lblNextFrame}{Next frame.}
\newcommand{\lblYourTurn}{Now one for you to do.}
\newcommand{\lblRouteTo}{to}


% ---------- Part titles ----------
% These live here rather than in structure.tex, and the reason is a rule this
% repository states and then broke once: nothing outside a language file may
% hard-code a word that the other edition also sets. Part titles written with
% \ifpl put nine user-visible strings outside the reach of the parity check
% that compares the two catalogues -- and, because \ifpl was fragile, produced
% a build failure in exactly one of the two editions.
\newcommand{\lblPartI}{Foundation}
\newcommand{\lblPartII}{Number, precision and cost}
\newcommand{\lblPartIII}{Linear algebra}
\newcommand{\lblPartIV}{Discrete structures and argument}
\newcommand{\lblPartV}{Calculus and automatic differentiation}
\newcommand{\lblPartVI}{Optimisation}
\newcommand{\lblPartVII}{Probability and statistics}
\newcommand{\lblPartVIII}{Information theory}
\newcommand{\lblPartIX}{Assembling it}
\newcommand{\lblPartAppendices}{Appendices}

% ---------- Instructions the method depends on ----------
\newcommand{\lblQuizIntro}{%
  Take this before you read the program and again after you have finished it.
  If you score well now, work only the frames the questions send you to. The
  frame numbers beside each question are the route back.}
\newcommand{\lblCanYouIntro}{%
  Rate yourself against each item, 1 (no) to 5 (yes, unaided). Anything below
  4 is not a matter for reflection: go back to the frames named beside it and
  work them again. The Summary's bracketed numbers are the same route.}
\newcommand{\lblCanYouFooter}{%
  Then take the Quiz a second time. It was a diagnostic on the way in and it
  is a test on the way out; the difference between the two scores is the only
  honest measure of what this program did for you.}
\newcommand{\lblTestIntro}{%
  These test exactly what the program taught, in the order it taught it.
  There are no traps here and nothing needs an idea you have not met.}
\newcommand{\lblFurtherIntro}{%
  Consolidation. Work these until the method is automatic rather than
  recalled --- that is the whole of the difference between having understood a
  technique and being able to use it.}
\newcommand{\lblAnswersIntro}{%
  Answers to the quizzes (Q), the test exercises (T) and the further problems
  (P) of every program, in program order.}

% The punctuation dash. English sets an em dash closed up; Polish sets a
% polpauza with a space on each side. Authors write \dash and never type a
% dash directly, because a translator cannot be expected to carry a typographic
% convention in their head forty-six programs deep.
\newcommand{\dash}{---}

% ---------- Admonition titles ----------
\newcommand{\lblNote}{Note}
\newcommand{\lblWarning}{Watch out}
\newcommand{\lblTrap}{The trap}
\newcommand{\lblInAI}{Where this shows up in AI}
\newcommand{\lblRigour}{What we are not proving}
\newcommand{\lblNotation}{Notation}
\newcommand{\lblVerify}{Check this before you trust it}
\newcommand{\lblExercise}{Exercise}

% ---------- Build-state markers ----------
\newcommand{\lblNotWritten}{NOT YET WRITTEN}
\newcommand{\lblNotRendered}{NOT RENDERED --- run \texttt{make diagrams}}
\newcommand{\lblSourceMissing}{Mermaid source missing}
\newcommand{\lblNoOutcomes}{No learning outcomes were declared for this program.}
\newcommand{\lblNoAnswers}{No answers were stored. This is a partial build.}
\newcommand{\lblDiagramManifest}{Diagram manifest}

% ============================================================
%  NOTATION CONTRACT — English edition
%  Every macro here has a counterpart in lang/pl.tex. Where the two editions
%  set different symbols, the difference is deliberate and is recorded in the
%  notation appendix.
% ============================================================
\newcommand{\mfadecimalmarker}{.}

% Trigonometric names. English uses tan/cot/arctan; Polish uses tg/ctg/arc tg,
% which is not a stylistic preference there but the universal convention.
\DeclareMathOperator{\tg}{tan}
\DeclareMathOperator{\ctg}{cot}
\DeclareMathOperator{\arctg}{arctan}
\DeclareMathOperator{\arcctg}{arccot}

% Probability. Both editions use Var and E, matching the literature the reader
% will actually be reading. The Polish edition carries a notation box saying
% so, because Polish textbooks classically write D^2(X) and M(X).
\newcommand{\Prob}{\mathbb{P}}
\newcommand{\Ex}{\mathbb{E}}
\DeclareMathOperator{\Var}{Var}
\DeclareMathOperator{\Cov}{Cov}
\DeclareMathOperator{\Corr}{Corr}

% Number theory. Polish writes NWD and NWW.
\DeclareMathOperator{\gcdop}{gcd}
\DeclareMathOperator{\lcmop}{lcm}

% Intervals. Polish tradition writes a closed interval with angle brackets,
% \langle a,b \rangle. Both editions use square brackets, because every paper
% the reader will meet does; the Polish edition says so in a notation box.
\newcommand{\intcc}[2]{[#1,\,#2]}
\newcommand{\intoo}[2]{(#1,\,#2)}
\newcommand{\intco}[2]{[#1,\,#2)}
\newcommand{\intoc}[2]{(#1,\,#2]}

% The date the book's facts were last checked. User-visible, so it lives here.
\newcommand{\pinnedon}{August 2026}
